% ===========================================================================
%  1 — Como ler este documento
% ===========================================================================
\chapter{Como ler este documento}
\label{cap:como-ler}

A Plataforma Unooria cresceu depressa: hoje ela tem nove missões para
estudantes, três ferramentas para professores, uma área inteira de gestão da
rede e um painel para a equipe da Unooria. Tudo isso foi construído, testado e
documentado do ponto de vista técnico --- mas nunca tinha sido descrito do
ponto de vista de \emph{produto}: quem usa, para quê, e o que conta como
\enquote{pronto}.

Este documento preenche essa lacuna. Ele é o \textbf{retrato do produto em
25 de setembro de 2026}, levantado a partir das telas, das regras de negócio e
dos testes de ponta a ponta que a plataforma já passa. Serve para três coisas:

\begin{enumerate}
  \item \textbf{Dar a todos a mesma visão do produto} --- o que ele é, para
    quem existe e quais promessas ele faz (Capítulo~\ref{cap:produto}).
  \item \textbf{Ser o backlog de partida} --- cada história tem estado, e o que
    está pela metade ou por fazer está consolidado no
    Capítulo~\ref{cap:backlog}.
  \item \textbf{Registrar as perguntas que só o produto pode responder} ---
    o Capítulo~\ref{cap:decisoes} lista as decisões em aberto encontradas no
    levantamento.
\end{enumerate}

\section{A hierarquia: épico, feature, história}

\begin{tabularx}{\linewidth}{@{}L{24mm}X L{30mm}@{}}
\toprule
\headrow \thd{Nível} & \thd{O que é} & \thd{Identificador}\\
\midrule
\textbf{Épico} & Um grande objetivo do produto, que atravessa várias telas e
  papéis. Ex.: \enquote{Cadastrar pessoas e turmas}. & \id{E01} a \id{E14}\\
\zebra \textbf{Feature} & Uma capacidade concreta dentro do épico, que alguém
  consegue usar e reconhecer. Ex.: \enquote{Importar pessoas por planilha}. &
  \id{F03.1} (épico 03, feature 1)\\
\textbf{História} & Uma necessidade de uma pessoa, contada do ponto de vista
  dela, com critérios que dizem quando está atendida. & \id{H03.1.2} (feature
  F03.1, história 2)\\
\bottomrule
\end{tabularx}

\section{O formato de cada história}

Toda história segue o mesmo molde, e o molde é o mesmo para o que já existe e
para o que falta:

\begin{historia}{H00.0.0}{Título curto da necessidade}{\pronto}
\hu{\emph{quem} (o papel, com o nome que a tela usa)}{\emph{o quê} (a ação, em
palavras de quem usa)}{\emph{por quê} (o valor para essa pessoa ou para a rede)}
\aceite
\begin{itemize}
  \item Cada critério é algo que dá para \textbf{verificar olhando a tela}.
  \item Os textos entre aspas são os que a interface mostra hoje, palavra por
    palavra.
\end{itemize}
\pendencia{quando algo da história ainda falta, aparece aqui.}
\end{historia}

Nas histórias marcadas \afazer, os critérios de aceite são uma
\textbf{proposta} para o PO validar, não uma decisão tomada --- o rótulo deixa
isso explícito.

\section{Os três estados}

\begin{tabularx}{\linewidth}{@{}L{24mm}X@{}}
\toprule
\headrow \thd{Estado} & \thd{Significado}\\
\midrule
\pronto & Funciona de ponta a ponta hoje, e há um teste automatizado que
  percorre a tela como uma pessoa faria.\\
\zebra \parcial & Funciona, mas falta uma parte declarada: um material
  (áudio, foto), um ajuste de tela, ou uma regra ainda divergente. A
  pendência está escrita na própria história.\\
\afazer & Ainda não existe. Está no documento porque apareceu como lacuna,
  risco ou evolução planejada.\\
\zebra \arquivada & Existiu e saiu do produto das escolas, mas ficou GUARDADA:
  outra missão ocupou o lugar dela na trilha, e ela não abre para ninguém. Só a
  equipe da Unooria vê a lista, e os dados das turmas continuam guardados.\\
\bottomrule
\end{tabularx}

\section{O que este documento \emph{não} é}

\begin{itemize}
  \item \textbf{Não é especificação técnica.} Não há aqui nomes de tabela,
    endereços internos nem desenho de sistema. Quem precisar desse nível
    encontra tudo na documentação técnica (PU-DOC-001); o
    Apêndice~\ref{ap:rastreio} diz em que capítulo dela cada épico mora.
  \item \textbf{Não é plano de lançamento.} A ordem dos épicos segue o caminho
    de uso da plataforma, não prioridade. A prioridade está no
    Capítulo~\ref{cap:backlog}.
  \item \textbf{Não é definitivo.} É a versão 0.1: o retrato do que existe,
    feito para ser corrigido pelo PO e mantido junto com o produto.
\end{itemize}

\begin{nota}
\textbf{Como manter vivo.} Uma história nova entra com o próximo número da
feature; uma história que muda de estado troca o selo. Nenhum identificador é
reaproveitado: se uma história for descartada, ela sai do documento e o número
fica vago. Assim, um \id{H04.2.3} citado num card ou numa conversa sempre aponta
para a mesma coisa.
\end{nota}
